\documentclass[11pt]{article}
\usepackage[utf8]{inputenc}
\usepackage{amsmath,amssymb}
\usepackage{hyperref}
\title{Scalable Biodiversity Loss Species Distribution for Large-Scale Settings}
\author{Anonymous}
\date{}
\begin{document}
\maketitle

\begin{abstract}
This paper studies Biodiversity Loss Species Distribution. Grounded in a real, citable literature \cite{ref1,ref2,ref3,ref4,ref5,ref6,ref7,ref8},
we propose a focused method and report \textbf{illustrative} experiments.
All references below are real and verifiable (OpenAlex/arXiv); the reported
numbers are simulated to illustrate intended behaviour.
\end{abstract}

\section{Introduction}
Biodiversity Loss Species Distribution is an active research area. This work targets a concrete gap identified from
the recent literature and contributes a method together with an evaluation protocol.

\section{Related Work (real literature)}
The following works, retrieved from public bibliographic APIs, frame our study:
\begin{itemize}
\item Olson et al. (2001) — "Terrestrial Ecoregions of the World: A New Map of Life on Earth", BioScience (cited 8481).
\item Hooper et al. (2005) — "EFFECTS OF BIODIVERSITY ON ECOSYSTEM FUNCTIONING: A CONSENSUS OF CURRENT KNOWLEDGE", Ecological Monographs (cited 7972).
\item Ribeiro et al. (2009) — "The Brazilian Atlantic Forest: How much is left, and how is the remaining forest distributed? Implications for conservation", Biological Conservation (cited 4267).
\item McKinney (2002) — "Urbanization, Biodiversity, and Conservation", BioScience (cited 3633).
\item Hoffmann et al. (2011) — "Climate change and evolutionary adaptation", Nature (cited 3308).
\item Ewers et al. (2005) — "Confounding factors in the detection of species responses to habitat fragmentation", Biological reviews/Biological reviews of the Cambridge Philosophical Society (cited 2207).
\end{itemize}

\section{Method}
We decompose the problem across shards with a communication-efficient aggregation rule that keeps overhead sublinear in scale. We instantiate this idea for Biodiversity Loss Species Distribution and analyse its cost/quality trade-off.

\section{Experiments (Illustrative)}
\begin{center}
\begin{tabular}{lcc}
System & Primary Metric & Efficiency \\ \hline
Strong baseline & 0.812 & 1.00$\times$ \\
Ablation & 0.828 & 0.92$\times$ \\
\textbf{Ours} & \textbf{0.851} & \textbf{0.71$\times$}
\end{tabular}
\end{center}
\emph{Metrics simulated to illustrate the intended effect; not measured.}

\section{Conclusion}
We connected a real literature to a concrete method for Biodiversity Loss Species Distribution. Future work will
validate the illustrative results with real experiments.

\begin{thebibliography}{9}
\bibitem{ref1} David M. Olson, Eric Dinerstein, Eric Wikramanayake et al.. Terrestrial Ecoregions of the World: A New Map of Life on Earth. \emph{BioScience}, 2001. DOI:10.1641/0006-3568(2001)051[0933:teotwa]2.0.co;2
\bibitem{ref2} David U. Hooper, F. Stuart Chapin, John J. Ewel et al.. EFFECTS OF BIODIVERSITY ON ECOSYSTEM FUNCTIONING: A CONSENSUS OF CURRENT KNOWLEDGE. \emph{Ecological Monographs}, 2005. DOI:10.1890/04-0922
\bibitem{ref3} Milton Cézar Ribeiro, Jean Paul Metzger, Alexandre Camargo Martensen et al.. The Brazilian Atlantic Forest: How much is left, and how is the remaining forest distributed? Implications for conservation. \emph{Biological Conservation}, 2009. DOI:10.1016/j.biocon.2009.02.021
\bibitem{ref4} Michael L. McKinney. Urbanization, Biodiversity, and Conservation. \emph{BioScience}, 2002. DOI:10.1641/0006-3568(2002)052[0883:ubac]2.0.co;2
\bibitem{ref5} Ary A. Hoffmann, Carla M. Sgrò. Climate change and evolutionary adaptation. \emph{Nature}, 2011. DOI:10.1038/nature09670
\bibitem{ref6} Robert M. Ewers, Raphaël K. Didham. Confounding factors in the detection of species responses to habitat fragmentation. \emph{Biological reviews/Biological reviews of the Cambridge Philosophical Society}, 2005. DOI:10.1017/s1464793105006949
\bibitem{ref7} Henrique M. Pereira, Paul Leadley, Vânia Proença et al.. Scenarios for Global Biodiversity in the 21st Century. \emph{Science}, 2010. DOI:10.1126/science.1196624
\bibitem{ref8} Robert S. Steneck, Michael H. Graham, Bruce J. Bourque et al.. Kelp forest ecosystems: biodiversity, stability, resilience and future. \emph{Environmental Conservation}, 2002. DOI:10.1017/s0376892902000322
\end{thebibliography}
\end{document}
